\subsection*{Two-sided physical source capacity}
Fix a true flight $j$ of an $m$-flight record, a marked unvisited disk,
and a clearance level $u_0\in[0,3/200]$. The original source consists
of two labelled circles of line directions $\vartheta$, one for each
impact sign. Its two halves contain exactly $j+1$ backward and
$m-j$ forward collisions from the marked foot. They are determined
by the same physical line; no independence or added reflection is used.
The level zero is the actual positive-side unvisited trace.

Theorem~\ref{thm:mb-branch-sup} strengthens the inherited momentum
kernel to a sum of suprema on disjoint genuine source intervals. If
$A_l=|\partial_\vartheta p_l|$ and $\mathcal P_l$ is its physical
half-word partition, then
\[
 \sum_{I\in\mathcal P_l}\sup_I A_l^{-1}\le H_{\rm nm},\qquad
 \sup_I A_l^{-1}\le E_*\sigma^{l-1},\qquad
 \sigma=47/59,\quad E_*=94/\sqrt{1827}.
\]
The first inequality uses the actual first-image density, the native
homogeneous growth bound, and the pointwise derivative at the final
true collision. It is not inferred from a density estimate at one
common terminal momentum. The second follows from the positive
physical optical product, including its exact endpoint convention.

The weighted interval-overlay lemma then pays the complete two-sided
source. With $u_x=|\partial_{\alpha_0}p_m|^{-1}$,
Theorem~\ref{thm:mb-two-sided-sum} proves
\begin{equation}\label{eq:intro-marked-capacity}
 \sum_{O\in\mathcal O_{m,j,R,u_0}}\sup_O u_x
 \le Q_{\rm mb}\bigl(\sigma^j+\sigma^{m-1-j}\bigr),\qquad
 Q_{\rm mb}=\frac{607}{90}H_{\rm nm}\frac{94}{\sqrt{1827}}.
\end{equation}
The sum contains both signs, every actual half-word intersection,
and arbitrarily short components. Different components of the same
real word are retained. The interval argument supplies the comparison
without a word-count factor. Original Borel record and roof selectors
are retained on these intervals, without charging their components again.

\subsection*{All critical offsets of the regular marked boundary}
The exact inherited boundary identity is
\[
 f^\varepsilon|d\theta|
 =\frac{W^\varepsilon}{4\pi Rc_*}
          \left|\partial_tu\big|_\theta\right|d\vartheta,
 \qquad \theta=\arg(-p_0,p_m).
\]
At a nonzero critical offset the physical normal speed can concentrate
near $\nabla F=0$. Its pointwise pole must be retained. The new
argument instead controls total turning on every actual overlap.
The absolute one-step logarithmic derivative, integrated on each
homogeneous image and summed by the adapted expansion, gives a
collision-uniform bound for $\operatorname{Var}\log A_l$.
This is an absolute-variation estimate, not an inference from a
H\"older difference bound.

For a regular plane curve $P$, radial turning is at most tangent
turning plus $\pi$. Allowing its one possible origin crossing costs
a second $\pi$ on the punctured intervals. Applying this exact
lemma to $P=(-p_0,p_m)$ gives
\[
 \int_{P\ne0}|d\arg P|\le 2\pi+V_A
\]
on every overlap, uniformly in its length and distance from criticality.
The original endpoint cap follows from the unchanged section.
The exact cofactor identity bounds
$f^\varepsilon\le\Lambda_\circ u_x/(4\pi c_*R_-^2)$.
Combining with \eqref{eq:intro-marked-capacity},
Theorem~\ref{thm:ab-all-offset} proves the complete regular marked
boundary estimate
\begin{equation}\label{eq:intro-all-offset}
 \mathfrak B^{\rm reg}_{j,u_0}
 \le C_{\rm ab}\bigl(\sigma^j+\sigma^{m-1-j}\bigr).
\end{equation}
There is no formal-domain coverage assumption, critical offset
restriction, lower incidence cutoff, or lower bound for $|\nabla F|$.
All actual regular pieces and the prescribed one-sided trace are summed
before the uniform estimate. In particular
\[
 \sum_{\min(j,m-1-j)\ge D}\mathfrak B^{\rm reg}_{j,u_j}
 \le\frac{2C_{\rm ab}}{1-\sigma}\sigma^D.
\]
For fixed proportional depth this full positive boundary charge tends
to zero after multiplication by $m^2$. Its pushforward by the original
roof has the same total-variation bound. This closes the nonzero-offset
\emph{regular angular-boundary} summation left by the preceding
zero-offset argument. It does not identify total variation with
essential height in physical time, or remove a distinct critical fan.

\subsection*{All-class physical outgoing-angle capacity}
The cross-curve theorem is also converted into a complete geometric
sum. The original outgoing-angle map
$\pi_R(s,\varphi)=s/R+\varphi\pmod{2\pi}$ has physical stable
affine fibers of exact length $\pi\sqrt{1+R^2}$. Treating their
first genuine backward collision before a current-strip subdivision
gives a uniform inverse-growth bound from each short affine test.
A fixed number of such tests covers every fiber, including its
collision-coordinate seam.

Let $\mathfrak C_N(R)$ consist of all actual backward word classes
with the precise homogeneity labels at $-1,\ldots,-N$, with every
disconnected component of one class counted together. Assigned
artificial strip faces remain in the partition. Set
\[
 d_N(C)=\sup_{y\in C\cap M_R'}j_N(T_R^{-N}y),
 \qquad \operatorname{cap}_R(C)=|\pi_R(C)|.
\]
Theorem~\ref{tc:all-classes} proves, pointwise in the outgoing
direction and uniformly in the original radius and collision count,
\[
 \sum_Cd_N(C)\one_{\pi_R(C)}(\vartheta)\le K_{\rm tc},
 \qquad \sum_Cd_N(C)\operatorname{cap}_R(C)\le2\pi K_{\rm tc}.
\]
The actual regular derivative at a strip-boundary representative is
paid by an adjacent true growth interval, with multiplicity at most
three. No bulk-null argument deletes those supremum values.

In particular the entire unweighted sum over classes containing an
actual unstable cone arc of fixed length $\delta>0$ is at most
$2\pi K_{\rm tc}/(\beta_{\rm tc}\delta)$.
This establishes a genuine long-class payment without a planar
continuation assumption. All remaining small-capacity classes are
still present in the exact pressure decomposition
\eqref{tc:remaining-pressure-interface}. Their full word-sup sum
requires an additional physical payment. Geometric outgoing-angle
capacity is not the original conditioned same-roof pool.

\subsection*{Positive critical fans even when the clearance slice collapses}
A distinct leading incoming critical measure remains when the regular
angular integral is punctured at zero gradient. A simultaneous
unvisited corner can have a positive physical fan while the fixed
zero-clearance angular slice has no regular interval.
Section~\ref{pc:section} pays this source on its actual physical side.

For every finite family of merged critical marked states, the finite
analytic physical conditions provide representatives at every common
sufficiently small positive clearance. A positive visited incidence
at each limiting collision makes its optical product continuous.
The strict momentum derivative prevents two distinct critical states
from occupying the same genuine one-sided source interval for all
small clearance. Thus both half-ray supremum sums from 319 transfer
to the critical family without requiring the zero-level slice to persist.
The exact two-half product now yields the stronger estimate
\[
 \mathfrak F_{m,j}(R)\le C_{\rm pc}
          \sigma^{\max(j,m-1-j)}.
\]
Here $\mathfrak F_{m,j}$ is the full original leading incoming
zero-clearance fan coefficient, with united one-copy physical
angular fractions and one original conditioning factor.
All strict words, both signs, and simultaneous original section
and unvisited cuts are retained. Theorem~\ref{pc:fan} gives
\[
 \sup_R m^2\sum_{j=0}^{m-1}\mathfrak F_{m,j}(R)
       \le C_{\rm pc}m^3\sigma^{(m-1)/2}\longrightarrow0.
\]
This includes every shallow marked index. No local-radius limit is
interchanged with the collision limit: finite-family approximation
proves an inequality at each fixed radius and collision count, whose
final constant is uniform. Actual visited-grazing traces and the
full collision-uniform local coarea remainders remain separate.
The full bulk singular-multiplier results of 315 stay inherited.

\subsection*{The inherited mathematical architecture}
All complete native proof bodies 260--318 are retained byte for byte
and included exactly once. Their results are inherited, including
the optical Jacobi and endpoint trace completion, exact marked
clearance flux and two-mark arithmetic, cross-flight and momentum
coarea, the native physical curve-growth theorem, complete regular
critical-sector coefficients, endpoint and radius resolvents,
physical tubes, and the full variational comparison chain.

In particular the global disk-product minimum is not assumed to be
a physical orbit. The retained common-gap and anchor-layer arguments
give the actual localization and clearance conclusions. With the
notation of Theorem~\ref{thm:ae-layer}, their energy estimate is
\begin{equation}\label{eq:ae-intro}
 \mathcal E_{d+94}\le q_{\rm a}\mathcal E_d,\qquad
 \mathcal E_d\le E_R(Q)q_{\rm a}^{\lfloor d/94\rfloor},\qquad
 q_{\rm a}=3263680/3263681.
\end{equation}
The exact physical clock keeps its endpoint work; the central record
transport keeps every original occupation test and separates the
terminal membership test. The subsequent full-source weak-roof,
singular-multiplier and regular-interior variation theorems remain
available with precisely their proved topologies. The common-word
cross-curve theorem retains its terminal incidence and physical
stable/optical factors. These completed results are not reintroduced
as new assumptions or new contributions.

The separately buildable historical main and technical supplement
retain the earlier geometric family, periodic orbits, arithmetic
obstructions, induced return, spectral, conditional-measure and
four-coordinate inversion architecture. Their original conditional
statements and correction records remain intact. The current
proof/dependency ledger distinguishes them from the unconditional
geometric deductions proved here.
